% !TEX root = root.tex
% !TEX recipe = LuaLaTeX+se
%\documentclass{article}
\documentclass[paper=letter,parskip]{scrartcl} % set document class: manual at https://ctan.org/pkg/koma-script
%\usepackage{fontspec}

% Load packages:
\usepackage[osf,p]{libertine} % set font
\usepackage[autocompile]{gregoriotex} % enable Gregorio score inclusion
\usepackage[latin]{babel} % set language
\usepackage{multicol}
\usepackage[margin=2cm]{geometry}
\usepackage{scrlayer-scrpage}

\setkomafont{section}{\normalfont\centering\huge\scshape} % section heading style
\setkomafont{subsection}{\normalfont\centering\Large\scshape\bfseries} % section heading style
\setkomafont{subsubsection}{\normalfont\centering\large\bfseries} % section heading style
\RedeclareSectionCommand[
    beforeskip=.5\baselineskip,
    runin=on]{paragraph}
\setcounter{secnumdepth}{-\maxdimen} % remove section numbering

\setkomafont{pagehead}{\scshape}
\automark[section]{section}
\automark*[subsection]{subsection}

%\lohead[\thesection]
%{\thesection}
%\rohead[\thesubsection]
%{\thesubsection}
%\pagestyle{scrheadings}


% Set the space around the initial:
% See http://gregorio-project.github.io/gregoriotex/details.html for more details and options
\grechangedim{beforeinitialshift}{2.2mm}{scalable}
\grechangedim{afterinitialshift}{2.2mm}{scalable}

% Set the initial font (change 43 for a larger size):
%\grechangestyle{initial}{\fontsize{43}{43}\selectfont}%

% Make staff lines red; remove for black:
%\gresetlinecolor{gregoriocolor}

% Use the "commentary" field of the score in the top right corner:
%\gresetheadercapture{commentary}{grecommentary}{string}

% Format annotation above initial
\grechangestyle{annotation}{\small\bfseries}


% DELTE THIS IF THE THING COMPILES
%\newcommand{\communio}[3]{
%    \greannotation{CO. #2}
%    \gregorioscore{#1/solesmes.gabc}
%    \IfFileExists{#1/english.gabc}{
%    \gregorioscore{#1/english.gabc}}{}
%    \gresetinitiallines{0}
%    \grecommentary{#3}
%    \gregorioscore{#1/verses.gabc}
%    \gresetinitiallines{1}
%}


\newcommand{\chnm}{undefined}

\begin{document}
\pagenumbering{arabic}

% \section{Foreword}
% \begin{itemize}
%     \item When a verse score annotation contains a star, (e.g., Ps. 37*,)
%         this indicates that another Psalm or canticle may be
%         chanted instead, likewise set to the same tone.
%     \item \emph{T.P.} means ``Tempus Paschale'' or Eastertide, typically
%         indicating an optional ``Alleluia'' that is to be sung only during
%         the Easter season (from Easter Sunday to Pentecost).
% \end{itemize}

%\clearpage
%% !TEX root = root.tex
% !TEX recipe = LuaLaTeX+se

\section{Tempus Adventus}

\subsection{Hebdomada I Adventus}
\greannotation{CO. I}
\gregorioscore{dominus-dabit/co--dominus_dabit--solesmes.1.gabc}
\gregorioscore{dominus-dabit/english.gabc}
\gresetinitiallines{0}
\gregorioscore{dominus-dabit/verses.gabc}
\gresetinitiallines{1}

\subsection{Hebdomada II Adventus}
\greannotation{CO. II}
\gregorioscore{jerusalem-surge/solesmes.gabc}
%\gregorioscore{jerusalem-surge/english.gabc}
\gresetinitiallines{0}
\grecommentary{Ps. 147}
\gregorioscore{jerusalem-surge/verses.gabc}
\gresetinitiallines{1}


\subsection{Hebdomada III Adventus}
\greannotation{CO. VII}
\gregorioscore{dicite-pusillanimes/solesmes.gabc}
%\gregorioscore{dicite-pusillanimes/english.gabc}
\gresetinitiallines{0}
\grecommentary{Is. 35}
\gregorioscore{dicite-pusillanimes/verses.gabc}
\gresetinitiallines{1}

\subsection{Hebdomada IV Adventus}
\begin{center}
\emph{T.P.: This chant may also be used for certain feasts or votive Masses
of the B.V.M.; when such a Mass falls during the Easter season, the Alleluia
is included.}
\end{center}
\greannotation{CO. I}
\grecommentary{Is. 7:14}
\gregorioscore{ecce-virgo/solesmes.gabc}
%\gregorioscore{ecce-virgo/english.gabc}
\gresetinitiallines{0}
\grecommentary{Ps. 18}
\gregorioscore{ecce-virgo/verses.gabc}
\gresetinitiallines{1}

\subsection{From December 17th to 24th}
\paragraph{Monday}~

\greannotation{CO. VI}
\grecommentary{Cf. Zach. 14: 5, 7}
\gregorioscore{ecce-dominus-veniet/solesmes.gabc}
% Ps. 49*, 1. 2-3 a. 3 be. 4· 5· 6

\paragraph{Tuesday}~

\greannotation{CO. VI}
\grecommentary{Ps. 18: 6, 7}
\gregorioscore{exsultavit-ut-gigas/solesmes.gabc}
% Ps. 18, 2. 3· 4· 5· 6 ab (Differentia : g)

\paragraph{Wednesday}~

\greannotation{CO. I}
\grecommentary{Cf. Hab. 3:3}
\gregorioscore{veni-domine-et-noli/solesmes.gabc}
% Ps. 84*, 2. 3· 4· 5· 7· 8. 10. 11. 12. 14

\paragraph{Thursday}
\emph{As on Monday.}
\paragraph{Friday}
\emph{As on Tuesday.}
\paragraph{Saturday}
\emph{As on Wednesday.}


%\clearpage
%% !TEX root = root.tex
% !TEX recipe = LuaLaTeX+se

\section{Tempus Nativitatis}
\subsection{In Nativitate Domini}
\subsubsection{At Vigil Mass}~

\greannotation{CO. I}
\grecommentary{Cf. Is. 40:5}
\gregorioscore{revelabitur/solesmes.gabc}
%\gregorioscore{revelabitur/english.gabc}
\gresetinitiallines{0}
\grecommentary{Ps. 23}
% 1. 2. 3. 4. 5. 6. 7. 8.
\gregorioscore{revelabitur/verses.gabc}
\gresetinitiallines{1}


\subsubsection{At Mass at Night}
\greannotation{CO. VI}
\grecommentary{Ps. 109:3}
\gregorioscore{in-splendoribus/solesmes.gabc}
%\gregorioscore{in-splendoribus/english.gabc}
\gresetinitiallines{0}
\grecommentary{Ps. 109}
% 1a. 1b. 2. 3. 4. 5. 7.
\gregorioscore{in-splendoribus/verses.gabc}
\gresetinitiallines{1}


\subsubsection{At Mass at Dawn}
\greannotation{CO. IV}
\grecommentary{Zach. 9:9}
\gregorioscore{exsulta-filia/solesmes.gabc}
%\gregorioscore{exsulta-filia/english.gabc}
\gresetinitiallines{0}
\grecommentary{Ps. 33*}
% It just says 33* (differentia g) ?
\gregorioscore{exsulta-filia/verses.gabc}\label{co:exsulta-filia}
\gresetinitiallines{1}


\subsubsection{At Mass at Daytime}
\greannotation{CO. I}
\grecommentary{Ps. 97:3cd}
\gregorioscore{viderunt-omnes/solesmes.gabc}
%\gregorioscore{viderunt-omnes/english.gabc}
\gresetinitiallines{0}
\grecommentary{Ps. 97}
% Ps. 97, I ab. I cd. 2. 3 ab. 4· 5· 6. 7· 8-9 a. 9 be
\gregorioscore{viderunt-omnes/verses.gabc}
\gresetinitiallines{1}


\subsection{Within the Octave of the Nativity}
\emph{All weekdays as on the Nativity at ``Mass at Dawn''
 or ``Mass at Daytime'', except:}
\paragraph{December 29th}
CO. Responsum,~\ref{co:responsum}.

\subsection{Holy Family of Jesus, Mary \& Joseph}
Celebrated on the Sunday within the Octave of the Nativity,
\emph{or} December 30, if the Nativity was on a Sunday.
\paragraph{Year A:}~

\greannotation{CO. VII}
\grecommentary{Mt. 2:20}
\gregorioscore{tolle-puerum/solesmes.gabc}
% %\gregorioscore{tolle-puerum/english.gabc}
% \gresetinitiallines{0}
% \grecommentary{Ps. 92*}
% %Ps. 92*, 1 ab. 1 c-2 a. 3· 4· 5
% %vel ps. 127*, 1. 2. 3· 4· 5· 6
% \gregorioscore{tolle-puerum/verses.gabc}
% \gresetinitiallines{1}

\paragraph{Years B \& C:}~

\greannotation{CO. I}
\grecommentary{Lk. 2:48, 49}
\gregorioscore{fili-quid/solesmes.gabc}
%\gregorioscore{fili-quid/english.gabc}
\gresetinitiallines{0}
\grecommentary{Ps. 26*}
%Ps. 26*, 1 a. 4 abc. 4 de. 5. 8. 13
\gregorioscore{fili-quid/verses.gabc}
\gresetinitiallines{1}


\subsection{Solemnity of Mary, Mother of God}
\emph{N.B.: The antiphon is as on the Nativity, at ``Mass at Dawn'',
but the verses are taken from the Magnificat.}
\greannotation{CO. IV}
\grecommentary{Zach. 9:9}
\gregorioscore{exsulta-filia/solesmes.gabc}
%\gregorioscore{exsulta-filia/english.gabc}
\gresetinitiallines{0}
\grecommentary{Luke 1:46-55}
\gregorioscore{exsulta-filia/magnificat-verses.gabc}
\gresetinitiallines{1}


\subsection{The Epiphany of the Lord}
\greannotation{CO. IV}
\grecommentary{Cf. Mt. 2:2}
\gregorioscore{vidimus-stellam/solesmes.gabc}
%\gregorioscore{vidimus-stellam/english.gabc}
\gresetinitiallines{0}
\grecommentary{Ps. 71*}
%Ps. 71*, 1. 2. 3. 7. 8. 10. 11. 12. 17 ab. 17 cd. 18
\gregorioscore{vidimus-stellam/verses.gabc}
\gresetinitiallines{1}

\paragraph{Weekdays after Epiphany.} As on Epiphany, or on the Nativity
of the Lord at Mass at Dawn.

\paragraph{January 7, where Epiphany is to be celebrated on Sunday, January 8th.}
CO. Dicit Dominus: Implete hydrias,~\ref{co:dicit-dominus-implete}.

\paragraph{January 8, or Tuesday after Epiphany}
CO. Manducaverunt,~\ref{co:Manducaverunt}.

\subsection{The Baptism of the Lord}
Ordinarily celebrated on the Sunday after Epiphany.
When the solemnity of the Epiphany of the Lord is transferred to a
Sunday falling on January 7 or 8, the feast of the Baptism of the Lord is
celebrated on the following Monday.
%\footnote{It is truly
%remarkable the lengths we will go to avoid just one more obligatory
%Mass during the year.}
\greannotation{CO. II}
\grecommentary{Gal. 3:27}
\gregorioscore{omnes-qui-in-christo/solesmes.gabc}
%\gregorioscore{omnes-qui-in-christo/english.gabc}
\gresetinitiallines{0}
\grecommentary{Ps. 28*}
%Ps. 28*, 1. 2. 3. 4. 5. 7-8. 10. 11
\gregorioscore{omnes-qui-in-christo/verses.gabc}
\gresetinitiallines{1}


%\clearpage
%% !TEX root = root.tex
% !TEX recipe = LuaLaTeX+se

%\section{Tempus Quadragesimae}

\subsection{First Sunday in Lent}
%\subsubsection{Dominica}
\renewcommand{\chnm}{scapulis-suis}
\greannotation{CO. III}
\grecommentary{Ps. 90:4-5}
\gregorioscore{\chnm/solesmes.gabc}
\IfFileExists{\chnm/english.gabc}{
\gregorioscore{\chnm/english.gabc}}{}
\gresetinitiallines{0}
\grecommentary{Ps. 90}
\gregorioscore{\chnm/verses.gabc}
\gresetinitiallines{1}
%Ps. 90, 1. 2. 3· 11. 12. 13. 14. 15.16

\clearpage
\subsection{Second Sunday in Lent}
%\subsubsection{Dominica}
\renewcommand{\chnm}{visionem}
\greannotation{CO. I}
\grecommentary{Mt. 17:9}
\gregorioscore{\chnm/solesmes.gabc}
\IfFileExists{\chnm/english.gabc}{
\gregorioscore{\chnm/english.gabc}}{}
\gresetinitiallines{0}
\grecommentary{Ps. 44}
\gregorioscore{\chnm/verses.gabc}
\gresetinitiallines{1}

\clearpage
\subsection{Third Sunday in Lent - Gospel of the Samaritan woman, option 1}
%\subsubsection{Dominica}
%If the Gospel about the Samaritan woman is read,
%use either of the following two.
\renewcommand{\chnm}{qui-biberit-3}
\greannotation{CO. III}
\grecommentary{John 4:13,14}
\gregorioscore{\chnm/solesmes.gabc}
\IfFileExists{\chnm/english.gabc}{
\gregorioscore{\chnm/english.gabc}}{}
\gresetinitiallines{0}
\grecommentary{Isaiah 12}
\gregorioscore{\chnm/verses.gabc}
\gresetinitiallines{1}

\clearpage
\subsection{Third Sunday in Lent - Gospel of the Samaritan woman, option 2}
%\subsubsection{Dominica}
%If the Gospel about the Samaritan woman is read,
%use either of the following two.
\renewcommand{\chnm}{qui-biberit-7}
\greannotation{CO. VII}
\grecommentary{John 4:13,14}
\gregorioscore{\chnm/solesmes.gabc}
\IfFileExists{\chnm/english.gabc}{
\gregorioscore{\chnm/english.gabc}}{}
\gresetinitiallines{0}
\grecommentary{Isaiah 12}
\gregorioscore{\chnm/verses.gabc}
\gresetinitiallines{1}


\clearpage
\subsection{Third Sunday in Lent - Other Gospel Readings}
\renewcommand{\chnm}{passer-invenit}
\greannotation{CO. I-A}
\grecommentary{Ps. 83:4-5}
\gregorioscore{\chnm/solesmes.gabc}
\IfFileExists{\chnm/english.gabc}{
\gregorioscore{\chnm/english.gabc}}{}
\gresetinitiallines{0}
\grecommentary{Ps. 83}
\gregorioscore{\chnm/verses.gabc}
\gresetinitiallines{1}


\clearpage
\subsection{Fourth Sunday in Lent - Gospel of the Man Born Blind}
\renewcommand{\chnm}{lutum-fecit}
\greannotation{CO. VI}
\grecommentary{John 9: 6,11,38}
\gregorioscore{\chnm/solesmes.gabc}
\IfFileExists{\chnm/english.gabc}{
\gregorioscore{\chnm/english.gabc}}{}
\gresetinitiallines{0}
\grecommentary{Ps. 26}
\gregorioscore{\chnm/verses.gabc}
\gresetinitiallines{1}
%Ps. 26*, 1 a. 1 b. 4 abc. 4 de. 5. 6 ab. 6 cd. 9 ab. 9 cd. 10. 13

\clearpage
\subsection{Fourth Sunday in Lent - Gospel of the Prodigal Son}
\renewcommand{\chnm}{oportet-te}
\greannotation{CO. VIII}
\grecommentary{Luke 15:32}
\gregorioscore{\chnm/solesmes.gabc}
\IfFileExists{\chnm/english.gabc}{
\gregorioscore{\chnm/english.gabc}}{}
\gresetinitiallines{0}
\grecommentary{Ps. 31}
\gregorioscore{\chnm/verses.gabc}
\gresetinitiallines{1}

\clearpage
\subsection{Fourth Sunday in Lent - Other Gospel Readings}
\renewcommand{\chnm}{jerusalem-quae-aedificatur}
\greannotation{CO. IV}
\grecommentary{Ps. 121:3-4}
\gregorioscore{\chnm/solesmes.gabc}
\IfFileExists{\chnm/english.gabc}{
\gregorioscore{\chnm/english.gabc}}{}
\gresetinitiallines{0}
\grecommentary{Ps. 121}
\gregorioscore{\chnm/verses.gabc}
\gresetinitiallines{1}

\clearpage
\subsection{Fifth Sunday in Lent - Gospel of Lazarus}
\renewcommand{\chnm}{videns-dominus}
\greannotation{CO. I}
\grecommentary{John. 11:33,35,43,44,39}
\gregorioscore{\chnm/solesmes.gabc}
\IfFileExists{\chnm/english.gabc}{
\gregorioscore{\chnm/english.gabc}}{}
\gresetinitiallines{0}
\grecommentary{Ps. 129}
\gregorioscore{\chnm/verses.gabc}
\gresetinitiallines{1}

\clearpage
\subsection{Fifth Sunday in Lent - Gospel of the Woman Caught in Adultery}
\renewcommand{\chnm}{nemo-te}
\greannotation{CO. VIII}
\grecommentary{John. 8:10-11}
\gregorioscore{\chnm/solesmes.gabc}
\IfFileExists{\chnm/english.gabc}{
\gregorioscore{\chnm/english.gabc}}{}
\gresetinitiallines{0}
\grecommentary{Ps. 31}
\gregorioscore{\chnm/verses.gabc}
\gresetinitiallines{1}

\clearpage
\subsection{Fifth Sunday in Lent - Other Gospel Readings}
\renewcommand{\chnm}{qui-mihi-ministrat}
\greannotation{CO. V}
\grecommentary{John. 12:26}
\gregorioscore{\chnm/solesmes.gabc}
\IfFileExists{\chnm/english.gabc}{
\gregorioscore{\chnm/english.gabc}}{}
\gresetinitiallines{0}
\grecommentary{Ps. 16} %I ab. I cd. 2. 3· 5· 6. 7· 8-9 a
%ps. 16*, 1 ab. 1 cd. 2. 9b-10. 11. 12. 15
\gregorioscore{\chnm/verses-lent.gabc}
\gresetinitiallines{1}




\clearpage
% !TEX root = root.tex
% !TEX recipe = LuaLaTeX+se

%\section{Tempus Quadragesimae}

\subsection{Palm Sunday}
\renewcommand{\chnm}{pater-si-non-potest}
\greannotation{CO. VIII}
\grecommentary{Mt. 26:42}
\gregorioscore{\chnm/solesmes.gabc}
\IfFileExists{\chnm/english.gabc}{
\gregorioscore{\chnm/english.gabc}}{}
\gresetinitiallines{0}
\grecommentary{Ps. 21} %I ab. I cd. 2. 3· 5· 6. 7· 8-9 a
%ps. 16*, 1 ab. 1 cd. 2. 9b-10. 11. 12. 15
\gregorioscore{\chnm/verses-21.gabc}
\gresetinitiallines{1}

\clearpage
\subsection{Maundy Thursday: Mass of the Lord's Supper}
\renewcommand{\chnm}{hoc-corpus}
\greannotation{CO. VIII}
\grecommentary{1 Cor. 11:24-25}
\gregorioscore{\chnm/solesmes.gabc}
\IfFileExists{\chnm/english.gabc}{
\gregorioscore{\chnm/english.gabc}}{}
\gresetinitiallines{0}
\emph{Choose from either Psalm 22 or Psalm 115 below.}
\subsubsection{Option One}
\grecommentary{Ps. 22} %I ab. I cd. 2. 3· 5· 6. 7· 8-9 a
%ps. 16*, 1 ab. 1 cd. 2. 9b-10. 11. 12. 15
\gregorioscore{\chnm/verses-22.gabc}

\subsubsection{Option Two}
\grecommentary{Ps. 115} %I ab. I cd. 2. 3· 5· 6. 7· 8-9 a
%ps. 16*, 1 ab. 1 cd. 2. 9b-10. 11. 12. 15
\gregorioscore{\chnm/verses-115.gabc}
\gresetinitiallines{1}





\clearpage
% !TEX root = root.tex
% !TEX recipe = LuaLaTeX+se

%\section{Tempus Quadragesimae}

\subsection{Sunday of the Resurrection - Vigil and Daytime Mass}
\renewcommand{\chnm}{pascha-nostrum}
\greannotation{CO. VI}
\grecommentary{1 Cor. 5:7-8}
\gregorioscore{\chnm/solesmes.gabc}
\IfFileExists{\chnm/english.gabc}{
\gregorioscore{\chnm/english.gabc}}{}
\gresetinitiallines{0}
\grecommentary{Ps. 117} %Ps. 117*, I. 2. 5· 8. 10. II. 13. 14. 15. 16. 17. 21. 22. 23. 24. 25.
%26. 28. 29
\gregorioscore{\chnm/verses.gabc}
\gresetinitiallines{1}

\clearpage
\subsection{Second Sunday of Easter}
\renewcommand{\chnm}{mitte-manum-alleluia}
\greannotation{CO. VI}
\grecommentary{John 20:27}
\gregorioscore{\chnm/solesmes.gabc}
\IfFileExists{\chnm/english.gabc}{
\gregorioscore{\chnm/english.gabc}}{}
\gresetinitiallines{0}
\grecommentary{Ps. 117} %Ps. 117*, I. 2. 5· 8. 10. II. 13. 14. 15. 16. 17. 21. 22. 23. 24. 25.
%26. 28. 29
\gregorioscore{pascha-nostrum/verses.gabc}
\gresetinitiallines{1}

\clearpage
\subsection{Third Sunday of Easter, Year A}
\renewcommand{\chnm}{surrexit-dominus}
\greannotation{CO. VI}
\grecommentary{Lk. 24:34}
\gregorioscore{\chnm/solesmes.gabc}
\IfFileExists{\chnm/english.gabc}{
\gregorioscore{\chnm/english.gabc}}{}
\gresetinitiallines{0}
\grecommentary{Ps. 117}
\gregorioscore{pascha-nostrum/verses.gabc}
\gresetinitiallines{1}



\clearpage
\subsection{Third Sunday of Easter, Year B}
\renewcommand{\chnm}{cantate-domino}
\greannotation{CO. II-A}
\grecommentary{Ps 95:2}
\gregorioscore{\chnm/solesmes.gabc}
\IfFileExists{\chnm/english.gabc}{
\gregorioscore{\chnm/english.gabc}}{}
\gresetinitiallines{0}
\grecommentary{Ps. 95}
\gregorioscore{\chnm/verses.gabc}
\gresetinitiallines{1}


\clearpage
\subsection{Third Sunday of Easter, Year C}
\renewcommand{\chnm}{simon-ioannis}
\greannotation{CO. VI}
\grecommentary{John 21:15,17}
\gregorioscore{\chnm/solesmes.gabc}
\IfFileExists{\chnm/english.gabc}{
\gregorioscore{\chnm/english.gabc}}{}
\gresetinitiallines{0}
\grecommentary{Ps. 33}
\gregorioscore{\chnm/verses-33.gabc}
\gresetinitiallines{1}

\clearpage
\subsection{Fourth Sunday of Easter} % p 224
\renewcommand{\chnm}{ego-sum-pastor} % with Alleluias
\greannotation{CO. II}
\grecommentary{John 10:14}
\gregorioscore{\chnm/solesmes-cum-alleluia.gabc}
\IfFileExists{\chnm/english.gabc}{
\gregorioscore{\chnm/english.gabc}}{}
\gresetinitiallines{0}
%Ps. 22*, I -2 a. 2 b-3 a. 3 b-4 ab. 4 cd. 5 ab. 5 cd. 6 ab. 6 cd.
%vel ps. 32*, 1. 12- 5. 18-22
\grecommentary{Ps. 22}
\gregorioscore{\chnm/verses.gabc}
\gresetinitiallines{1}


\clearpage
\subsection{Fifth Sunday of Easter, Year A} % p 560
\renewcommand{\chnm}{tanto-tempore} % with Alleluias
\greannotation{CO. IV}
\grecommentary{John 14:9}
\gregorioscore{\chnm/solesmes.gabc}
\IfFileExists{\chnm/english.gabc}{
\gregorioscore{\chnm/english.gabc}}{}
\gresetinitiallines{0}
% cum ps. 32*, I. 2. 3· 12. 13. 18
\grecommentary{Ps. 32}
\gregorioscore{\chnm/verses.gabc}
\gresetinitiallines{1}

\clearpage
\subsection{Fifth Sunday of Easter (Other years)} % p 228
\renewcommand{\chnm}{ego-sum-vitis-vera} % with Alleluias
\greannotation{CO. VIII}
\grecommentary{John 15:5}
\gregorioscore{\chnm/solesmes.gabc}
\IfFileExists{\chnm/english.gabc}{
\gregorioscore{\chnm/english.gabc}}{}
\gresetinitiallines{0}
%Ps. 79*, 2 ab. 9· 10. 11. 12. 16. 18. 19
\grecommentary{Ps. 79}
\gregorioscore{\chnm/verses.gabc}
\gresetinitiallines{1}



\clearpage
\subsection{Sixth Sunday of Easter, Year A} % p 232
\renewcommand{\chnm}{non-vos-relinquam} % with Alleluias
\greannotation{CO. V}
\grecommentary{John 14:18}
\gregorioscore{\chnm/solesmes.gabc}
\IfFileExists{\chnm/english.gabc}{
\gregorioscore{\chnm/english.gabc}}{}
\gresetinitiallines{0}
%Ps. 121*, 1. 2. 3· 4· 5· 6. 7· 8. 9
\grecommentary{Ps. 121}
\gregorioscore{\chnm/verses.gabc}
\gresetinitiallines{1}


\clearpage
\subsection{Sixth Sunday of Easter, Year B} % p 436
\renewcommand{\chnm}{ego-vos-elegi} % with Alleluias
\greannotation{CO. I}
\grecommentary{John 15:16}
\gregorioscore{\chnm/solesmes.gabc}
\IfFileExists{\chnm/english.gabc}{
\gregorioscore{\chnm/english.gabc}}{}
\gresetinitiallines{0}
%Ps. 88*, 2. 4· 6. 20. 21. 22. 25. 29 (Differentia : a)
\grecommentary{Ps. 88}
\gregorioscore{\chnm/verses.gabc}
\gresetinitiallines{1}

\clearpage
\subsection{Sixth Sunday of Easter, Year C} % p 232
\renewcommand{\chnm}{spiritus-sanctus-docebit} % with Alleluias
\greannotation{CO. VIII}
\grecommentary{John 14:26}
\gregorioscore{\chnm/solesmes.gabc}
\IfFileExists{\chnm/english.gabc}{
\gregorioscore{\chnm/english.gabc}}{}
\gresetinitiallines{0}
% Ps. 50*, 3 a. 9· 10. 12. 13. 15. 17. 20 (Differentia : G*)
\grecommentary{Ps. 50}
\gregorioscore{\chnm/verses.gabc}
\gresetinitiallines{1}


\clearpage
\subsection{The Ascension of the Lord, Year A} % p 213
\renewcommand{\chnm}{data-est-mihi} % with Alleluias 
\greannotation{CO. I}
\grecommentary{Mt. 28: 18, 19}
\gregorioscore{\chnm/solesmes.gabc}
\IfFileExists{\chnm/english.gabc}{
\gregorioscore{\chnm/english.gabc}}{}
\gresetinitiallines{0}
% cum ps. 77*, 1. 3-4 a. 23. 24. 25. 27
\grecommentary{Ps. 77}
\gregorioscore{\chnm/verses.gabc}
\gresetinitiallines{1}

\clearpage
\subsection{The Ascension of the Lord, Year B} % p 437
\renewcommand{\chnm}{signa-eos} % with Alleluias Me. 16, 17. 18
\greannotation{CO. VII}
\grecommentary{Mc. 16: 17, 18}
\gregorioscore{\chnm/solesmes.gabc}
\IfFileExists{\chnm/english.gabc}{
\gregorioscore{\chnm/english.gabc}}{}
\gresetinitiallines{0}
% cum ps. 33*.
\grecommentary{Ps. 33}
\gregorioscore{\chnm/verses.gabc}
\gresetinitiallines{1}


\clearpage
\subsection{The Ascension of the Lord, Year C} % p 238
\renewcommand{\chnm}{psallite-domino} % with Alleluias
\greannotation{CO. I}      %CHANGE
\grecommentary{Ps. 67:33-34}    %CHANGE
\gregorioscore{\chnm/solesmes.gabc}
\IfFileExists{\chnm/english.gabc}{
\gregorioscore{\chnm/english.gabc}}{}
\gresetinitiallines{0}
% Ps. 67, 2. 5 abc. 5 d- 6. 19. 20. 21. 25. 29. 30. 33-34
\grecommentary{Ps. 67}        %CHANGE
\gregorioscore{\chnm/verses.gabc}
\gresetinitiallines{1}


\clearpage
\subsection{Seventh Sunday of Easter} % p 243
\emph{This is not celebrated in most dioceses in the United States,
as, regrettably, the feast of the Ascension is celebrated on this Sunday
instead.}
\renewcommand{\chnm}{pater-cum-essem} % with Alleluias
\greannotation{CO. IV}      %CHANGE
\grecommentary{John 17: 12, 13, 15}    %CHANGE
\gregorioscore{\chnm/solesmes.gabc}
\IfFileExists{\chnm/english.gabc}{
\gregorioscore{\chnm/english.gabc}}{}
\gresetinitiallines{0}
% Ps. 121*, 1. 2. 3· 4· 5· 6. 7· 8. 9
\grecommentary{Ps. 121}        %CHANGE
\gregorioscore{jerusalem-quae-aedificatur/verses.gabc}
\gresetinitiallines{1}


\clearpage
\subsection{Pentecost Sunday, at the Vigil Mass} % p 251
\renewcommand{\chnm}{ultimo-festivitatis} % with Alleluias
\greannotation{CO. V}      %CHANGE
\grecommentary{John 7:37-39}    %CHANGE
\gregorioscore{\chnm/solesmes.gabc}
\IfFileExists{\chnm/english.gabc}{
\gregorioscore{\chnm/english.gabc}}{}
\gresetinitiallines{0}
%Ps. 103*, 1 ab. 30. 31. 33· 34
\grecommentary{Ps. 103}        %CHANGE
\gregorioscore{\chnm/verses.gabc}
\gresetinitiallines{1}

\clearpage
\subsection{Pentecost Sunday, at the Daytime Mass} % p 256
\renewcommand{\chnm}{factus-est-repente} % with Alleluias
\greannotation{CO. VII}      %CHANGE
\grecommentary{Acts 2:2,4}    %CHANGE
\gregorioscore{\chnm/solesmes.gabc}
\IfFileExists{\chnm/english.gabc}{
\gregorioscore{\chnm/english.gabc}}{}
\gresetinitiallines{0}
%Ps. 67*, 2. 4· 5 abc. 5 d- 6. 8. 9· 20. 21. 29. 36
\grecommentary{Ps. 67}        %CHANGE
\gregorioscore{\chnm/verses.gabc}
\gresetinitiallines{1}



% \section{Proprium de Sanctis}
% \subsubsection{February 2, Presentation of the Lord}
% \label{co:responsum}

\end{document}